% Chapitre 2 — Contexte du projet
\chapter{Contexte du projet}
\label{ch:contexte}

\section{Présentation de l'organisme d'accueil}

\textbf{Talys Consulting} est un cabinet de conseil spécialisé dans la transformation digitale, l'intelligence artificielle et l'accompagnement des institutions financières. L'entreprise intervient sur des problématiques de :
\begin{itemize}[leftmargin=*]
  \item modélisation du risque crédit et conformité KYC/AML ;
  \item industrialisation de pipelines data et MLOps ;
  \item intégration de solutions IA dans les systèmes d'information existants (SI bancaires, plateformes ICM).
\end{itemize}

Le projet PFE s'inscrit dans une démarche d'innovation visant à prototyper une brique de scoring intelligente réutilisable dans des contextes clients microfinance et banque inclusive.

\section{Cadre académique — ESPRIT}

Ce projet est réalisé dans le cadre du \textbf{cycle ingénieur} à l'\textbf{École Supérieure Privée d'Ingénierie et de Technologies (ESPRIT)}. La formation en ingénierie informatique y combine :
\begin{itemize}[leftmargin=*]
  \item fondements théoriques (algorithmique, probabilités, bases de données) ;
  \item compétences en génie logiciel (UML, architecture, tests) ;
  \item projets applicatifs en entreprise (PFE obligatoire).
\end{itemize}

L'encadrement académique assure le respect des standards de rédaction scientifique, la cohérence méthodologique et l'évaluation finale du livrable.

\section{Encadrement}

\begin{table}[H]
  \centering
  \caption{Composition de l'équipe d'encadrement}
  \begin{tabular}{|l|l|p{7cm}|}
    \hline
    \textbf{Rôle} & \textbf{Nom} & \textbf{Responsabilités} \\
    \hline
    Encadrant entreprise & Oussema Midouni & Cadrage métier, revue architecture, priorités \\
    \hline
    Encadrant entreprise & Hamza Ksentini & Spécifications REF-08, validation dataset, revue technique \\
    \hline
    Encadrante académique & Salma Hajjem & Suivi PFE, relecture rapport, évaluation ESPRIT \\
    \hline
    Stagiaire ingénieur & Chedi Bouali & Conception, développement, tests, rédaction \\
    \hline
  \end{tabular}
\end{table}

\section{Planning du projet}

Le projet s'est déroulé sur une période estimée à \textbf{16 semaines}, réparties selon le planning suivant :

\begin{table}[H]
  \centering
  \caption{Planning macro du projet PFE}
  \label{tab:planning}
  \begin{tabular}{|c|l|p{8cm}|}
    \hline
    \textbf{S.} & \textbf{Phase} & \textbf{Activités principales} \\
    \hline
    1--2 & Cadrage & Lecture REF-08, analyse existant, setup environnement \\
    \hline
    3--4 & Données & Génération Faker, pipeline ETL, feature engineering \\
    \hline
    5--6 & Modèles classiques & LR, RF, XGBoost, métriques, API /predict \\
    \hline
    7--8 & KYC + LLM & Score KYC, Ollama, endpoint /explain \\
    \hline
    9--10 & Séquentiel & LSTM/GRU, transactions + remboursements \\
    \hline
    11--12 & Graphe & GraphSAGE, réseau relationnel \\
    \hline
    13--14 & Agent + RAG & LangGraph, TF-IDF, rapports Markdown/PDF \\
    \hline
    15--16 & UI + validation & React, tests, documentation, rapport PFE \\
    \hline
  \end{tabular}
\end{table}

\section{Document de référence REF-08}

Le document \textbf{REF-08 — Spécifications Techniques} constitue la source normative du projet. Ses exigences principales sont :

\subsection{Modèles d'intelligence artificielle}

\begin{itemize}[leftmargin=*]
  \item \textbf{Modèles classiques} : Régression Logistique, Forêt Aléatoire, XGBoost ;
  \item \textbf{Modèles séquentiels} : LSTM, GRU pour l'analyse des transactions ;
  \item \textbf{Modèles de graphes} : GNN / GraphSAGE pour l'analyse du réseau relationnel.
\end{itemize}

\subsection{Sorties des modèles}

Chaque analyse doit produire :
\begin{itemize}[leftmargin=*]
  \item \keyword{Score KYC} (0--100) ;
  \item \keyword{Score de risque crédit} / niveau (FAIBLE, MODERE, ELEVE) ;
  \item \keyword{Probabilité de défaut} ;
  \item \keyword{Explication du score} (Explainable AI).
\end{itemize}

\subsection{Agent conversationnel}

\begin{itemize}[leftmargin=*]
  \item Basé sur LangChain / LangGraph ;
  \item Tools pour appeler dynamiquement les modèles ML ;
  \item Mémoire conversationnelle ;
  \item Génération de rapports (LLM + RAG).
\end{itemize}

\section{Environnement technique}

\begin{table}[H]
  \centering
  \caption{Stack technologique du projet}
  \begin{tabular}{|l|l|}
    \hline
    \textbf{Couche} & \textbf{Technologies} \\
    \hline
    Langage & Python 3.11+, TypeScript \\
    \hline
    Data & Pandas, NumPy, Parquet, Faker \\
    \hline
    ML classique & scikit-learn, XGBoost, Joblib \\
    \hline
    Deep Learning & PyTorch (LSTM, GRU, GraphSAGE) \\
    \hline
    API & FastAPI, Pydantic, Uvicorn \\
    \hline
    LLM local & Ollama (llama3.2) \\
    \hline
    Orchestration & LangChain, LangGraph \\
    \hline
    RAG & TF-IDF (scikit-learn), index local Joblib \\
    \hline
    Frontend & React, Vite, TypeScript \\
    \hline
    Export & ReportLab (PDF), Markdown \\
    \hline
    IDE / OS & Cursor, Windows 10/11, Git \\
    \hline
  \end{tabular}
\end{table}
